\documentclass[10pt]{article}
\usepackage{titlesec}
\usepackage{geometry}
\geometry{verbose,tmargin=.9in,bmargin=.9in,lmargin=1.0in,rmargin=1.0in}
\usepackage{amsmath,amsfonts,amsthm,amssymb}
\usepackage{url}
\usepackage{color}
\usepackage[usenames,dvipsnames,svgnames,table]{xcolor}
\usepackage[colorlinks=true, linkcolor=red, urlcolor=blue, citecolor=gray]{hyperref}
\usepackage{float}
\usepackage{caption}
\usepackage{subcaption}
\usepackage{graphicx}
\usepackage{wrapfig}
\usepackage{booktabs}
\usepackage{longtable}
\usepackage{enumitem}
\usepackage{multicol}
\usepackage{etoolbox}
%\usepackage{bbm}
\usepackage{ulem}

\newcommand{\bs}[1]{\boldsymbol{#1}}
\newcommand{\bv}[1]{\mathbf{#1}}
\newcommand{\R}{\mathbb{R}}
\newcommand{\E}{\mathbb{E}}
\DeclareMathOperator*{\argmin}{arg\,min}
\DeclareMathOperator*{\argmax}{arg\,max}
\DeclareMathOperator{\rank}{rank}
\DeclareMathOperator{\Var}{Var}
\DeclareMathOperator{\Cost}{Cost}
\DeclareMathOperator{\cut}{cut}
\DeclareMathOperator{\tr}{tr}

\newcommand{\norm}[1]{\left\lVert #1\right\rVert}
\newcommand{\ip}[2]{\left\langle #1,\, #2\right\rangle}

\definecolor{nyuDarkPurple}{HTML}{330662}
\definecolor{nyuOfficialPurple}{HTML}{57068c}

\newcommand{\spara}[1]{\vspace{.5em}\noindent {\large\sffamily\textcolor{nyuOfficialPurple}{#1}}}
\titleformat{\section}[hang]{\Large\sffamily\color{nyuDarkPurple}}{\thesection}{1em}{}
\titleformat{\subsection}[hang]{\large\sffamily\color{nyuDarkPurple}}{\thesection}{1em}{}
\titleformat{\subsubsection}[hang]{\normalsize\sffamily\color{gray}}{\thesection}{1em}{}

\usepackage{fancyhdr}
\setlength{\headheight}{30pt}
\pagestyle{fancy}
\lhead{}
\rhead{\thepage}
\pagenumbering{gobble}

\setcounter{secnumdepth}{0}

\begin{document}
	
\begin{center}
	\normalsize
	New York University Tandon School of Engineering
	
	Computer Science and Engineering
	\medskip
	
	\large
	CS-GY 6763: Homework 4.
	
	\medskip
	
	\normalsize 
	\noindent \emph{Collaboration is allowed on this problem set, but solutions must be written-up individually. Please list collaborators for each problem separately, or write ``No Collaborators" if you worked alone.}
	\medskip
\end{center}

\subsection{Problem 1: Approximating Eigenvalues Moments}
\textbf{(15 pts)} Let $A \in \R^{n\times n}$ be a square symmetric matrix, which means it is guaranteed to have a symmetric eigendecomposition with real eigenvalues, $\lambda_1 \ge \ldots \ge \lambda_n$, and orthogonal eigenvectors. While computing these eigenvalues naively takes $O(n^3)$ time, we can compute their sum much more quickly: with $n$ operations. This is because $\sum_{i=1}^n \lambda_i$ is exactly equal to the trace of $A$, i.e. the sum of its diagonal entries $\tr(A) = \sum_{i=1}^n A_{ii}$. We can also compute the sum of squared eigenvalues in $O(n^2)$ time by taking advantage of the fact that $\sum_{i=1}^n \lambda_i^2 = \norm{A}_F^2$ where $\norm{A}_F^2$ is the Frobenius norm. What about $\sum_{i=1}^n \lambda_i^3$? Or $\sum_{i=1}^n \lambda_i^4$? It turns out that no exact algorithms faster than a full eigendecomposition are known. In this problem, however, we show how to approximate $\sum_{i=1}^n \lambda_i^k$ for any positive integer $k$ in $O(n^2 k)$ time. Doing so turns out to have a lot of applications in machine learning and data science.

\begin{enumerate}
	\item Show that $\sum_{i=1}^n \lambda_i^k = \tr(A^k)$ where $A^k$ denotes the chain of matrix products $A \cdot A \cdot \ldots \cdot A$, repeated $k$ times. For the remainder of the problem we use the notation $B = A^k$.
	\item Let $x^{(1)}, \ldots, x^{(m)} \in \R^n$ be $m$ independent random vectors, all with i.i.d. $\{+1, -1\}$ uniform random entries. Let $Z = \frac{1}{m} \sum_{i=1}^m (x^{(i)})^T B x^{(i)}$. We will show that $Z$ is a good estimator for $\tr(B)$ and thus for $\sum_{i=1}^n \lambda_i^k$. Give a short argument that $Z$ can be computed in $O(n^2 k m)$ time (recall that $B = A^k$).
	\item Prove that:
	\[
	\E[Z] = \tr(B)
	\qquad\text{and}\qquad
	\Var[Z] \le \frac{2}{m} \norm{B}_F^2.
	\]
	Hint: Use linearity of variance but be careful about what things are independent!
	\item Show that if $m = O(1/\epsilon^2)$ then, with probability $9/10$,
	\[
	|\tr(B) - Z| \le \epsilon \norm{B}_F.
	\]
	\item Argue that, when $A$ is positive semidefinite (equivalently, $\lambda_1, \ldots, \lambda_n$ are all positive), $\epsilon \norm{B}_F \le \epsilon \tr(B)$, so the above guarantee actually gives the relative error bound,
	\[
	(1-\epsilon)\tr(B) \le Z \le (1+\epsilon)\tr(B),
	\]
	all with just $O(n^2 k/\epsilon^2)$ computation time.
\end{enumerate}

\spara{Pr. 1(a)}
Let $A\in\R^{n\times n}$ be symmetric with eigenvalues $\lambda_1,\dots,\lambda_n$. Write the orthonormal eigendecomposition
\[
A = U\Lambda U^\top,
\qquad
\Lambda=\mathrm{diag}(\lambda_1,\dots,\lambda_n),
\qquad
U^\top U=I.
\]
Then
\[
A^k=(U\Lambda U^\top)^k = U\Lambda^k U^\top.
\]
Using cyclicity of trace,
\[
\tr(A^k)=\tr(U\Lambda^k U^\top)=\tr(\Lambda^k U^\top U)=\tr(\Lambda^k)
=\sum_{i=1}^n \lambda_i^k.
\]
Thus,
\[
\boxed{\sum_{i=1}^n \lambda_i^k=\tr(A^k).}
\]

\medskip
\spara{Pr. 1(b)}
Let $B:=A^k$. Let $x^{(1)},\dots,x^{(m)}\in\{\pm 1\}^n$ be independent random vectors with i.i.d.\ entries taking values $\pm 1$ with equal probability, and define
\[
Z=\frac{1}{m}\sum_{i=1}^m (x^{(i)})^\top Bx^{(i)}.
\]

\textbf{How to compute $Z$ }(I did this without forming $B$).
For each $i=1,\dots,m$, set $v^{(0)}:=x^{(i)}$ and iterate
\[
v^{(t)} := A v^{(t-1)} \quad \text{for }t=1,2,\dots,k.
\]
Then $v^{(k)} = A^k x^{(i)} = Bx^{(i)}$. Define
\[
s_i := (x^{(i)})^\top v^{(k)} = (x^{(i)})^\top Bx^{(i)},
\qquad
Z=\frac{1}{m}\sum_{i=1}^m s_i.
\]

\textbf{Runtime.}
Each dense matrix--vector multiply $Av$ costs $O(n^2)$. We do $k$ multiplies per sample, so each $s_i$ costs $O(kn^2)$, and therefore
\[
\boxed{\text{time to compute }Z\text{ is }O(mkn^2).}
\]

\medskip
\textbf{Trace connection (in expectation).}
For any single sample $x\in\{\pm 1\}^n$, define $s:=x^\top Bx$. Expanding,
\[
s = x^\top Bx = \sum_{p=1}^n\sum_{q=1}^n B_{pq}\,x_p x_q.
\]
Taking expectation and using linearity,
\[
\E[s] = \sum_{p=1}^n\sum_{q=1}^n B_{pq}\,\E[x_p x_q].
\]
Because the entries are i.i.d.\ and $\E[x_p]=0$ by definition, for $p\neq q$,
\[
\E[x_p x_q] = \E[x_p]\E[x_q] = 0,
\]
and for $p=q$,
\[
\E[x_p^2]=1.
\]
Therefore only diagonal terms remain:
\[
\E[s] = \sum_{p=1}^n B_{pp} = \tr(B).
\]
In particular, for $s_i=(x^{(i)})^\top Bx^{(i)}$,
\[
\E[s_i]=\tr(B)
\quad\Rightarrow\quad
\E[Z]=\E\left[\frac{1}{m}\sum_{i=1}^m s_i\right]
=\frac{1}{m}\sum_{i=1}^m \E[s_i]
=\tr(B).
\]

\medskip
\spara{Pr. 1(c)}
Let $Y:=x^\top Bx$ for a single random vector $x\in\{\pm 1\}^n$ with i.i.d.\ entries.
Then $Z=\frac{1}{m}\sum_{i=1}^m Y_i$ where $Y_i:=(x^{(i)})^\top Bx^{(i)}$ are i.i.d.

\textbf{Expectation.}
As shown above,
\[
\E[Y]=\tr(B),
\qquad\text{hence}\qquad
\E[Z]=\tr(B).
\]

\textbf{Variance reduction across $m$ samples.}
Since $Y_1,\dots,Y_m$ are independent and identically distributed,
\[
\Var(Z)
=
\Var\left(\frac{1}{m}\sum_{i=1}^m Y_i\right)
=
\frac{1}{m^2}\sum_{i=1}^m \Var(Y_i)
=
\frac{1}{m}\Var(Y).
\]

\textbf{Key identity (I got a bit stuck here until I rembered it):}
\[
\boxed{\Var(Y)=\E[Y^2]-\big(\E[Y]\big)^2.}
\]

\textbf{Compute $\E[Y^2]$.}
Write
\[
Y=x^\top Bx=\sum_{p,q} B_{pq}\,x_p x_q.
\]
Then
\[
Y^2
=
\left(\sum_{p,q} B_{pq}\,x_p x_q\right)
\left(\sum_{r,s} B_{rs}\,x_r x_s\right)
=
\sum_{p,q,r,s} B_{pq}B_{rs}\,x_p x_q x_r x_s.
\]
Taking expectation,
\[
\E[Y^2]=\sum_{p,q,r,s} B_{pq}B_{rs}\,\E[x_p x_q x_r x_s].
\]
Because the coordinates are i.i.d.\ with $\E[x_j]=0$ and $x_j^2=1$, the expectation
$\E[x_p x_q x_r x_s]$ is $0$ unless indices pair up (each index appears an even number of times).
The surviving pairings yield
\[
\E[Y^2]=(\tr(B))^2 + 2\|B\|_F^2.
\]
Therefore,
\[
\Var(Y)=\E[Y^2]-(\E[Y])^2
=\big((\tr(B))^2 + 2\|B\|_F^2\big)-(\tr(B))^2
=2\|B\|_F^2.
\]
Finally,
\[
\boxed{
\Var(Z)=\frac{1}{m}\Var(Y)\le \frac{2}{m}\|B\|_F^2.
}
\]

\medskip
\spara{Pr. 1(d)}

From part (c), we have
\[
\E[Z]=\tr(B),
\qquad
\Var(Z)\le \frac{2}{m}\|B\|_F^2.
\]
By Chebyshev's inequality, for any $t>0$,
\[
\Pr\big(|Z-\E[Z]|\ge t\big)\le \frac{\Var(Z)}{t^2}.
\]
Set $t=\varepsilon\|B\|_F$. Then
\begin{align*}
\Pr\big(|Z-\tr(B)|\ge \varepsilon\|B\|_F\big)
&=
\Pr\big(|Z-\E[Z]|\ge t\big) \\
&\le
\frac{\Var(Z)}{\varepsilon^2\|B\|_F^2}
\le
\frac{(2/m)\|B\|_F^2}{\varepsilon^2\|B\|_F^2}
=
\frac{2}{m\varepsilon^2}.
\end{align*}
Choose $m$ so that $\frac{2}{m\varepsilon^2}\le \frac{1}{10}$, i.e.
\[
m \ge \frac{20}{\varepsilon^2}.
\]
Equivalently, taking $m=\left\lceil \frac{20}{\varepsilon^2}\right\rceil$ yields
\[
\Pr\big(|Z-\tr(B)|\le \varepsilon\|B\|_F\big)\ge \frac{9}{10}.
\]
In particular,
\[
\boxed{m=O\!\left(\frac{1}{\varepsilon^2}\right).}
\]

\medskip
\spara{Pr. 1(e)}

Assume $A\succeq 0$ (PSD). Recall $B=A^k$.

From part (d) we already have, with probability at least $9/10$,
\[
|Z-\tr(B)|\le \varepsilon \|B\|_F.
\]
We will show that in the PSD case this implies the relative-error guarantee
\[
\Pr\big(|Z-\tr(B)|\le \varepsilon\,\tr(B)\big)\ge \frac{9}{10}.
\]

\medskip
\textbf{Key inequality: $\|B\|_F \le \tr(B)$ when $B\succeq 0$.}
If $A\succeq 0$, then all eigenvalues $\lambda_i\ge 0$. Hence $B=A^k\succeq 0$ and the eigenvalues of $B$ are $\lambda_i^k\ge 0$.
Therefore,
\[
\tr(B)=\sum_{i=1}^n \lambda_i^k,
\qquad
\|B\|_F=\sqrt{\sum_{i=1}^n (\lambda_i^k)^2}
=\sqrt{\sum_{i=1}^n \lambda_i^{2k}}.
\]
For nonnegative numbers $a_i\ge 0$ we have $\sqrt{\sum_i a_i^2}\le \sum_i a_i$.
Applying this with $a_i=\lambda_i^k$ gives
\[
\boxed{\ \|B\|_F \le \tr(B)\ }.
\]

\medskip
\textbf{Conclusion (absolute error $\Rightarrow$ relative error).}
Since $\|B\|_F \le \tr(B)$, we have $\varepsilon\|B\|_F \le \varepsilon\,\tr(B)$. Thus
\[
|Z-\tr(B)|\le \varepsilon\|B\|_F
\quad\Longrightarrow\quad
|Z-\tr(B)|\le \varepsilon\,\tr(B).
\]
Hence the same probability bound from part (d) applies:
\[
\Pr\big(|Z-\tr(B)|\le \varepsilon\,\tr(B)\big)\ \ge\ \Pr\big(|Z-\tr(B)|\le \varepsilon\|B\|_F\big)\ \ge\ \frac{9}{10}.
\qquad \Box
\]

\medskip
\textbf{Compute time (I did this niavely by forming $B=A^k$ directly ... I think there may be a shortcut gut feeling though).}
Alternatively, we can explicitly compute the matrix $B=A^k$ first, and then evaluate the quadratic forms.
Assume $A$ is dense $n\times n$.

\begin{itemize}
\item \emph{Forming $B=A^k$.} Compute $A^2,A^3,\dots,A^k$ by repeated matrix--matrix multiplication:
\[
A^2=A\cdot A,\quad A^3=A^2\cdot A,\ \dots,\ A^k=A^{k-1}\cdot A.
\]
Each dense matrix--matrix multiply costs $O(n^3)$, and we perform $(k-1)$ such multiplies, so forming $B$ costs
\[
O(kn^3).
\]

\item \emph{Evaluating the $m$ samples.} For each sample $x^{(i)}$, compute $y^{(i)}=Bx^{(i)}$ in $O(n^2)$ time (dense matrix--vector multiply), then compute $(x^{(i)})^\top y^{(i)}$ in $O(n)$ time. Thus each sample costs $O(n^2)$, and $m$ samples cost
\[
O(mn^2).
\]
\end{itemize}

Therefore, the total running time using this approach is
\[
\boxed{O(kn^3 + mn^2).}
\]

\subsection{Problem 2: Gradient Descent as Power Method}
\textbf{(10 pts)} Consider minimizing the function $f(x) = \norm{Ax-b}_2^2$ for $A \in \R^{n\times d}$ and $b \in \R^n$ with gradient descent. $f(x)$ has gradient $\nabla f(x) = 2A^T A x - 2A^T b$. As mentioned in class, we can obtain a convergence bound for this function that scales with the condition number of $A^T A$, which is the ratio of the largest and smallest eigenvalues of $A^T A$, $\lambda_1/\lambda_d$. It turns out that this can be proven using an analysis very similar to power method.

\begin{enumerate}
	\item Prove that, when gradient descent is run on $f(x)$, the difference between the $i$th iterate and $x^*$ can be written as:
	\[
	x^{(i)} - x^* = (I - 2\eta A^T A)(x^{(i-1)} - x^*)
	\]
	By induction, it follows that the error $x^{(i)} - x^*$ equals $(I - 2\eta A^T A)^i (x^{(0)} - x^*)$.
	\item Prove that if we set $\eta = 1/(2\lambda_1)$, then after $i = O\bigl(\tfrac{\lambda_1}{\lambda_d}\log(1/\epsilon)\bigr)$ iterations, $\norm{(I - \tfrac{1}{\lambda_1} A^T A)^i}_2 < \epsilon$. Recall that $\norm{\cdot}_2$ denotes the spectral norm of a matrix, which is equivalent to its largest singular value, or the largest absolute value of the matrix's eigenvalues.
	\item Conclude that, after $O\bigl(\tfrac{\lambda_1}{\lambda_d}\log(1/\epsilon)\bigr)$ iterations, gradient descent finds a vector $x^{(i)}$ with $\norm{x^{(i)} - x^*}_2 \le \epsilon\norm{x^{(0)} - x^*}_2$. This can be a one line proof.
	\item We can also use this approach to obtain an ``accelerated'' version of gradient descent. To do so, prove that, for any degree $q$ polynomial $p = c_0 + c_1 x + \ldots + c_q x^q$ with $p(1) = c_0 + c_1 + \ldots + c_q = 1$ and any starting vector $x^{(0)}$, we can compute using $q$ gradient computations and up to $O(dq)$ additional runtime a vector $x^{(q)}$ such that:
	\[
	x^{(q)} - x^* = p\Bigl(I - \tfrac{1}{\lambda_1} A^T A\Bigr)(x^{(0)} - x^*).
	\]
	Note that this result strictly generalizes what we know from gradient descent, which computes $x^{(q)}$ satisfying the equation for the polynomial $p(x) = x^q$, which satisfies our restriction that $p(1) = 1$.
	\item Prove that for $q = O\bigl(\tfrac{\lambda_1}{\lambda_d}\log(1/\epsilon)\bigr)$, there exists a polynomial $p$ with coefficients $c_0 + c_1 + \ldots + c_q = 1$ such that $\norm{p(I - \tfrac{1}{\lambda_1} A^T A)}_2 \le \epsilon$. Hint: You might want to use Claim 4 in the supplemental notes on the Lanczos method posted on the course website.
\end{enumerate}

Just as for regular gradient descent, it follows that $\norm{x^{(q)} - x^*}_2 = \norm{p(I - \tfrac{1}{\lambda_1} A^T A)(x^{(0)} - x^*)}_2 \le \epsilon\norm{x^{(0)} - x^*}_2$ as long as we use degree $q = O\bigl(\tfrac{\lambda_1}{\lambda_d}\log(1/\epsilon)\bigr)$ -- i.e. use $O\bigl(\tfrac{\lambda_1}{\lambda_d}\log(1/\epsilon)\bigr)$ gradient computations.

\subsection{Problem 3: Spectral Methods for Cliques}
\textbf{(15 pts)} A common task in data mining is to identify large cliques in a graph. For example, in social networks, large cliques can be indicators of fraudulent accounts or networks of accounts designed to promote certain content. In this problem, we consider a spectral heuristic for finding a large clique based on the top eigenvector of the graph adjacency matrix $A$:
\begin{itemize}
	\item Compute the leading eigenvector $v_1$ of $A$.
	\item Let $i_1, \ldots, i_k \in \{1, \ldots, n\}$ be the indices of the $k$ entries in $v_1$ with largest absolute value.
	\item Check if nodes $i_1, \ldots, i_k$ form a $k$-clique.
\end{itemize}

We will analyze this heuristic on a natural random graph model. Specifically, let $G$ be an Erdos--Renyi random graph: we start with $n$ nodes, and for every pair of nodes $(i, j)$, we add an edge between the pair with probability $p < 1$. To simplify the math, also assume that we add a self-loop at every vertex $i$ with probability $p$. Then, choose a fixed subset $S$ of $k$ nodes to form a clique. Connect all nodes in $S$ with edges and add self-loops. We will argue that, for sufficiently large $k$, we can expect the heuristic above to identify the nodes in the clique.

\begin{enumerate}
	\item Let $A$ be the adjacency matrix of a random graph generated as above. What is $\E[A]$? Prove that the rank of $\E[A]$ is $2$. In other words, the matrix only has two non-zero eigenvalues.
	\item Derive expressions for the two non-zero eigenvalues of $\E[A]$, and their corresponding eigenvectors. Hint: First argue that, up to multiplying by a constant, any eigenvector $v$ must have $v[i] = 1$ for all $i \notin S$ and $v[i] = \alpha$ for all $i \in S$, where $\alpha$ is a constant. Then use some high school algebra.
	\item Using your results from (2) above, argue that, up to a positive scaling, the top eigenvector $v_1$ has $v[i] = 1$ for all $i \notin S$ and $v[i] = \alpha$ for all $i \in S$, where $\alpha > 1$. In other words, the largest entries of $v_1$ exactly correspond to the nodes in the clique.
	\item To prove the algorithm works, it is possible to use a matrix concentration inequality to argue that the top eigenvector of $A$ is close to that of $\E[A]$. Instead of doing that, let's verify things experimentally. Generate a graph $G$ according to the prescribed model with $n = 900$, $k = |S| = 30$, and $p = .1$. Compute the top eigenvector of $A$ and look at its 30 largest entries in magnitude. What fraction of nodes in the clique $S$ are among these 30 entries? Repeat the experiment and report the average fraction recovered.
\end{enumerate}

\spara{Pr. 3(1)}

\textbf{Step 1: Compute $\E[A]$ connectivity matric for each entry.}
Since $A_{ij}\in\{0,1\}$,
\[
(\E[A])_{ij}=\E[A_{ij}]=\Pr(A_{ij}=1).
\]
By the model:
\begin{itemize}
\item If $i\in S$ and $j\in S$, the edge is forced, so $\Pr(A_{ij}=1)=1$.
\item Otherwise, the edge is random with probability $p$, so $\Pr(A_{ij}=1)=p$.
\end{itemize}
Therefore,
\[
(\E[A])_{ij}=
\begin{cases}
1 & \text{if } i\in S \text{ and } j\in S,\\
p & \text{otherwise.}
\end{cases}
\]

\medskip
\textbf{Fact (rank via row/column space).}
\[
\boxed{
\rank(M)=\dim(\text{RowSpace}(M))=\dim(\text{ColSpace}(M)).
}
\]
In particular, if all rows of $M$ lie in the span of at most $r$ vectors, then $\rank(M)\le r$.

\medskip
\textbf{Step 2: Show $\rank(\E[A])\le 2$.}
Consider the rows of $\E[A]$. There are only two row types:
\begin{itemize}
\item If $i\notin S$, then row $i$ has $(\E[A])_{ij}=p$ for every $j$, i.e.\ the row is all $p$'s.
\item If $i\in S$, then row $i$ has $(\E[A])_{ij}=1$ for $j\in S$ and $(\E[A])_{ij}=p$ for $j\notin S$.
\end{itemize}
Call these two row vectors $r_{\mathrm{out}}$ and $r_{\mathrm{in}}$. Every row of $\E[A]$ is either $r_{\mathrm{out}}$ or $r_{\mathrm{in}}$, hence every row lies in
\[
\mathrm{span}\{r_{\mathrm{out}},\,r_{\mathrm{in}}\}.
\]
By the rank fact above,
\[
\boxed{\rank(\E[A]) \le 2.}
\]

\spara{Pr. 3(2)}

Because $\rank(M)=2$ and the matrix treats all clique vertices the same and all non-clique vertices the same, eigenvectors live in the $2$D space ``constant on $S$'' and ``constant on $S^c$.'' Thus solving $Mv=\lambda v$ reduces to a $2\times 2$ system.

\medskip
\textbf{Goal for (2).}
Find the two nonzero eigenvalues/eigenvectors of $M=\E[A]$. We already know $\rank(M)=2$, so there are exactly $2$ nonzero eigenvalues.

From (1), entrywise:
\[
M_{ij}=
\begin{cases}
1 & i\in S,\ j\in S,\\
p & \text{otherwise.}
\end{cases}
\]
Let $|S|=k$ and $n-k$ be outside, and write $n_0:=n-k$.

\medskip
\textbf{Using the hint here : eigenvector has only two values.}
Let $v$ be of the form
\[
v_i=
\begin{cases}
\alpha & i\in S,\\
1 & i\notin S.
\end{cases}
\]
We want $Mv=\lambda v$. Because there are only two ``types'' of coordinates, we write the eigenvector equation for:
\begin{enumerate}
\item a typical coordinate outside the clique ($i\notin S$),
\item a typical coordinate inside the clique ($i\in S$).
\end{enumerate}

\medskip
\textbf{Case 1: $i\notin S$ (outside).}
Row $i$ of $M$ is all $p$'s, so
\[
(Mv)_i = \sum_{j=1}^n p\,v_j = p\sum_{j=1}^n v_j.
\]
But
\[
\sum_{j=1}^n v_j = n_0\cdot 1 + k\cdot \alpha = n_0+k\alpha,
\]
so
\[
(Mv)_i = p(n_0+k\alpha).
\]
Since $v_i=1$ outside, $(Mv)_i=\lambda v_i=\lambda$, hence
\[
\boxed{\lambda = p(n_0+k\alpha).}
\]

\medskip
\textbf{Case 2: $i\in S$ (inside).}
Row $i$ has entries $1$ on columns in $S$ and entries $p$ on columns not in $S$, so
\[
(Mv)_i
=
\sum_{j\in S} 1\cdot v_j + \sum_{j\notin S} p\cdot v_j
=
k\alpha + pn_0.
\]
Since $v_i=\alpha$ inside, $(Mv)_i=\lambda v_i=\lambda\alpha$, hence
\[
\boxed{k\alpha + pn_0 = \lambda\alpha.}
\]

\medskip
\textbf{Now it is a $2\times 2$ system.}
We can view this as two equations in the two unknowns $(\lambda,\alpha)$:
\[
\lambda = p(n_0+k\alpha),
\qquad
\lambda\alpha = k\alpha + pn_0.
\]

\medskip
\textbf{Solve for $\alpha$ (quadratic form).}
Substitute $\lambda=p(n_0+k\alpha)$ into $\lambda\alpha=k\alpha+pn_0$:
\[
\alpha\cdot p(n_0+k\alpha)=k\alpha+pn_0
\quad\Longleftrightarrow\quad
pk\,\alpha^2 + (pn_0-k)\alpha - pn_0 = 0.
\]
Let $\alpha_\pm$ be the two roots of this quadratic. Then the corresponding eigenvalues are
\[
\boxed{\lambda_{\pm}=p\big(n_0+k\alpha_{\pm}\big).}
\]

\medskip
\textbf{Write the eigenvectors.}
Each eigenvector has the two-level form:
\[
v^{(\pm)}_i=
\begin{cases}
\alpha_{\pm} & i\in S,\\
1 & i\notin S.
\end{cases}
\]
Therefore the two nonzero eigenpairs are $(\lambda_+,v^{(+)})$ and $(\lambda_-,v^{(-)})$.
Since $\rank(M)=2$, all other eigenvalues are $0$.

\spara{Pr. 3(3)}

\textbf{What we already have from (2).}
We consider eigenvectors of $M=\E[A]$ of the form
\[
v_i=
\begin{cases}
1 & i\notin S,\\
\alpha & i\in S,
\end{cases}
\]
and we derived the two equations (with $n_0=n-k$):
\[
\text{(outside)}\quad \lambda = p(n_0+k\alpha),
\qquad
\text{(inside)}\quad k\alpha+pn_0=\lambda\alpha.
\]
Eliminating $\lambda$ gives that $\alpha$ satisfies
\[
pk\alpha^2 + (pn_0-k)\alpha - pn_0 = 0.
\]
Let the two roots be $\alpha_+$ and $\alpha_-$. These correspond to the two nonzero eigenvectors/eigenvalues.

\medskip
\noindent
\textbf{Claim.} The eigenvector corresponding to the largest eigenvalue has $\alpha>1$ (in fact, $\alpha_+>1$).

\medskip
\textbf{Step 1: Largest eigenvalue $\Longleftrightarrow$ largest $\alpha$.}
From the outside equation,
\[
\lambda = p(n_0+k\alpha).
\]
Since $p>0$ and $n_0,k$ are fixed positive constants, $\lambda$ is an increasing linear function of $\alpha$. Therefore, the top eigenvalue corresponds to the larger root $\alpha_+$, and it suffices to show
\[
\alpha_+>1.
\]

\medskip
\textbf{Step 2: Evaluate the quadratic at $\alpha=1$.}
Define
\[
f(\alpha):=pk\alpha^2+(pn_0-k)\alpha-pn_0.
\]
The roots of $f$ are exactly $\alpha_-$ and $\alpha_+$. Evaluate at $\alpha=1$:
\[
f(1)=pk+(pn_0-k)-pn_0 = pk-k = k(p-1).
\]
Assuming $0<p<1$, we have $p-1<0$, so
\[
f(1)=k(p-1)<0.
\]

\medskip
\textbf{Step 3: Conclude $\alpha_+>1$.}
Since $pk>0$, the quadratic $f(\alpha)$ opens upward and
\[
f(\alpha)\to +\infty \quad \text{as }\alpha\to\infty.
\]
We already computed $f(1)<0$. Because $f$ is continuous and goes from negative at $\alpha=1$ to positive for large $\alpha$, there exists some $\alpha>1$ with $f(\alpha)=0$. That root is the larger root $\alpha_+$, hence
\[
\boxed{\alpha_+>1.}
\]

\medskip
\textbf{Interpretation.}
The top eigenvector has larger entries on clique nodes than non-clique nodes (outside entries are $1$, inside entries are $\alpha_+>1$). Therefore the largest coordinates of the top eigenvector correspond to the clique. Intuitively this makes sense because the clique block has all $1$ connections between its members, whereas outside the clique edges appear only with probability $p<1$.

\spara{Pr. 3(4) (experimental verification)}
\textbf{Goal.}
We experimentally test whether the top eigenvector of the random adjacency matrix $A$ places clique
nodes among its largest coordinates.

\medskip
\textbf{Simulation (one trial).}
We generate a graph according to the planted-clique model with $n=900$, $k=|S|=30$, and $p=0.1$:
we sample a clique set $S\subseteq [n]$ of size $k$, set $A_{ij}=1$ for all $i,j\in S$
(clique block, including diagonal self-loops), and for all other pairs $(i,j)$ set $A_{ij}=1$
independently with probability $p$.
We compute the top eigenvector $v$ of $A$, take the $30$ indices with largest magnitude entries $|v_i|$,
call this set $T$, and report the recovery fraction
\[
\mathrm{frac}=\frac{|S\cap T|}{30}.
\]

\medskip
\textbf{Results.}
Over $20$ independent trials, the average fraction recovered was
\[
\boxed{\overline{\mathrm{frac}}\approx 0.885}
\]
(with standard deviation $\approx 0.048$, minimum $0.80$, and maximum $0.967$).
Equivalently, on average the method recovered about $0.885\times 30 \approx 26.6$ of the $30$ clique
nodes among the top-$30$ eigenvector entries.

\end{document}